\documentclass[11pt]{article}
\usepackage[margin=1in]{geometry}
\usepackage{amsmath,amssymb,amsthm}
\usepackage{graphicx}
\usepackage[hidelinks]{hyperref}

\newtheorem{theorem}{Theorem}
\newtheorem{lemma}{Lemma}
\newtheorem{remark}{Remark}

\newcommand{\B}{\mathcal B}
\newcommand{\Calk}{\mathcal C}
\newcommand{\Czero}{\mathcal C_0}
\newcommand{\Fzero}{F_0}
\newcommand{\comm}{\operatorname{comm}}

\title{A Necessary and Sufficient Condition for \(\pi(I-P+K)\) to be Generalized Drazin Invertible}
\author{agent\_lane\_04}
\date{June 29, 2026}

\begin{document}
\maketitle

\section*{Status}

This packet gives a likely valid full answer to the open question appearing
in the example before Section 3 of Hamdan--Berkani \cite{HB240101390}.  In
that example \(P\in L(X)\) is an idempotent, \(K\in L(X)\) is compact, and
\[
  T=I-P+K.
\]
The paper gives a sufficient commutant condition for \(\pi(T)\) to be
generalized Drazin invertible in
\(\Czero(X)=L(X)/\Fzero(X)\), then asks more generally for conditions on
\(P\) and \(K\).  The result below shows that the stated sufficient condition
is also necessary for this \(T=I-P+K\) setup.

\section*{Source Question}

The source passage occurs in the example immediately after Remark 2.13, around
line 847 of the cached source for arXiv:2401.01390.  In text form:
\begin{quote}
More generally, what are the conditions on \(P\) and \(K\) so that
\(\pi(T)\) is generalized Drazin invertible in \(\mathcal C_0(X)\)?
\end{quote}
The surrounding example has fixed \(T=I-P+K\), where \(P\) is idempotent and
\(K\) is compact.

\begin{center}
\includegraphics[width=0.86\textwidth]{figures/open_problem_crop.png}
\end{center}

\section*{Definitions and Source Input}

Let \(X\) be a Banach space.  Write
\[
  \pi:L(X)\to \Czero(X)=L(X)/\Fzero(X)
\]
for the quotient map by the ideal of finite-rank operators, and write
\[
  \Pi:L(X)\to \Calk(X)=L(X)/K(X)
\]
for the Calkin quotient by compact operators.  For an element \(a\) in an
algebra, \(\comm(a)\) denotes its commutant.

We use the following result from the source paper.

\begin{lemma}[Hamdan--Berkani, Remark 2.13]
\label{lem:source}
Let \(S\in L(X)\) be a pseudo B-Fredholm operator, and let \(P\) be an
idempotent associated to \(S\) in the Calkin algebra \(\Calk(X)\).  Then
\(\pi(S)\) is generalized Drazin invertible in \(\Czero(X)\) if and only if
\[
  \comm(\pi(S))\subseteq \comm(\pi(P)).
\]
\end{lemma}

In the source example with \(T=I-P+K\), \(K\) compact, the paper observes that
\(T\) is pseudo B-Fredholm and that \(P\) is associated to \(T\) in
\(\Calk(X)\).  Thus Lemma \ref{lem:source} applies to this \(T\).

\section*{Proof Intuition}

Only one algebraic identity is missing from the source's sufficient-condition
argument.  In \(\Czero(X)\), set
\[
  t=\pi(T), \qquad p=\pi(P), \qquad k=\pi(K).
\]
Since \(T=I-P+K\), we have \(t=1-p+k\), equivalently \(p=1+k-t\).  Therefore,
for every \(x\) commuting with \(t\),
\[
  xp-px = x(1+k-t)-(1+k-t)x = xk-kx.
\]
So, on the commutant of \(t\), commuting with \(p\) is exactly the same as
commuting with \(k\).

\section*{Main Theorem}

\begin{theorem}
\label{thm:main}
Let \(X\) be a Banach space, \(P\in L(X)\) an idempotent, \(K\in L(X)\) a
compact operator, and \(T=I-P+K\).  Put \(t=\pi(T)\) and \(k=\pi(K)\).  Then
\[
  t \text{ is generalized Drazin invertible in } \Czero(X)
  \quad\Longleftrightarrow\quad
  \comm(t)\subseteq \comm(k).
\]
Equivalently,
\[
  \begin{aligned}
  &\pi(I-P+K) \text{ is generalized Drazin invertible in } \Czero(X)\\
  &\qquad\Longleftrightarrow\quad
  \comm(\pi(I-P+K))\subseteq \comm(\pi(P)).
  \end{aligned}
\]
\end{theorem}

\begin{proof}
Let \(p=\pi(P)\).  The source example shows that \(T\) is pseudo B-Fredholm
and that \(P\) is associated to \(T\) in \(\Calk(X)\).  By Lemma
\ref{lem:source},
\[
  t \text{ is generalized Drazin invertible in } \Czero(X)
  \quad\Longleftrightarrow\quad
  \comm(t)\subseteq \comm(p).
\]

It remains only to compare the two commutant inclusions.  Since
\[
  t=\pi(I-P+K)=1-p+k,
\]
we have \(p=1+k-t\).  If \(x\in\comm(t)\), then
\[
  xp-px=x(1+k-t)-(1+k-t)x=xk-kx.
\]
Hence, for every \(x\in\comm(t)\), \(x\in\comm(p)\) if and only if
\(x\in\comm(k)\).  Therefore
\[
  \comm(t)\subseteq\comm(p)
  \quad\Longleftrightarrow\quad
  \comm(t)\subseteq\comm(k).
\]
Combining this equivalence with Lemma \ref{lem:source} proves the theorem.
\end{proof}

\section*{Consequences}

The sufficient condition stated in the source,
\[
  \comm(\pi(T))\subseteq \comm(\pi(K)),
\]
is not merely sufficient for \(T=I-P+K\); it is necessary and sufficient.  In
particular, the obstruction can be stated entirely in terms of \(P\) and \(K\):
every operator that commutes with \(I-P+K\) modulo finite rank must commute
with \(K\) modulo finite rank, equivalently with \(P\) modulo finite rank.

\section*{Verification Notes}

The proof has no computational component.  The only external mathematical
input is the source paper's Remark 2.13, together with the source example's
observation that \(T=I-P+K\) is pseudo B-Fredholm with \(P\) associated to
\(T\) in the Calkin algebra.  The added argument is the identity
\([x,\pi(P)]=[x,\pi(K)]\) for all \(x\in\comm(\pi(T))\).

\section*{Limitations}

This theorem answers the open question in the natural scope of the displayed
source example \(T=I-P+K\).  It does not attempt to classify arbitrary compact
perturbations of unrelated pseudo B-Fredholm decompositions beyond that
normal form.

\section*{Novelty Check}

The local run indexes were searched for arXiv:2401.01390, the source title,
``conditions on \(P\) and \(K\)'', ``\(I-P+K\)'', ``generalized Drazin
invertible'', and the commutant-inclusion phrasing.  No duplicate packet or
attempt was found.  A bounded web search on June 29, 2026 for the exact open
question language, the source title, and the phrase ``\(T=I-P+K\)'' with
``generalized Drazin invertible'' found no later explicit answer.

\section*{Human Review Recommendation}

Review the interpretation of the open question as referring to the immediately
preceding \(T=I-P+K\) setup.  Under that reading, the proof is complete and
the necessary direction is the only point not written in the source example.

\begin{thebibliography}{9}

\bibitem{HB240101390}
A. Hamdan and M. Berkani,
\emph{Fredholm-type Operators and Index},
arXiv:2401.01390, 2024.

\bibitem{Koliha}
J. J. Koliha,
\emph{A generalized Drazin inverse},
Glasgow Mathematical Journal 38 (1996), 367--381.

\end{thebibliography}

\end{document}
